\section{Interpretive Boundaries and Future Validation}\label{sec:interp}
\subsection{Why the target determines the conclusion}
The conservative analysis attempts to avoid treating a disputed augmented assertion as a true defect. The broader runner-only analysis instead asks whether an original process measure ranks \emph{whatever the existing runner reports as failure}. Either may be relevant to a different operational goal, but they are not interchangeable. For true \emph{semantic correctness}, neither label policy supplies a fully validated reference standard.

This distinction is not unique to the present study. SWE-bench Verified was introduced as a human-filtered subset of the original benchmark \cite{openai2024verified}; later commentary raised concerns about training-data exposure and benchmark task/test design \cite{openai2026benchmark}. This external critique does \emph{not} establish contamination of the Gemini 3 Flash source studied here; no sample-specific contamination test was conducted. SWE-Cycle evaluates environment setup, code generation and verification together, and proposes hybrid execution-aware judgments \cite{guan2026swecycle}. Such designs provide competing alternatives to relying on one runner outcome.

\subsection{Proposed independent semantic annotation extension}
A high-priority extension would independently assess the issue-level relevance of the 31 historical runner FAILs and relevant coverage-limited records. To prevent confirmation bias, experts should receive the original issue, exact model patch, pinned augmented-test diff, base environment and observed test output \emph{without seeing the trajectory score or desired study result}. At least two experienced engineers would independently record one of four categories per disputed assertion: (i) issue-required semantic defect; (ii) unsupported or wrong augmented assertion; (iii) environment/harness-confounded; (iv) unresolved from available evidence. An adjudicator could resolve disagreements without rewriting the frozen historical results.

Reviewer agreements, rubric versions and unresolved cases would be archived separately. A future analysis could then compare documentary labels with expert judgments, add distinct positive tasks, and repeat the study on another agent/model. This paragraph is a \emph{prospective proposal}, not a report of completed human validation.

\subsection{Why not report a calibrated detector today?}
The earlier synthetic paper used logistic regression and a six-feature composite with simulated outcomes. Neither has been re-fitted on real runs. The current five positive task families cannot justify stable deployment thresholds, externally assessed probability calibration or causal conclusions. PR-AUC and decision curves would become meaningful for a future deployment-oriented question only after label adjudication, adequate distinct positive tasks and a genuine holdout design. Their absence is a limitation for practical prediction, not a result that should be invented to satisfy an old synthetic-study reviewer checklist.

\subsection{Negative findings are not proof of independence}
The results do not prove that clean process and correct patches are unrelated. They show that four particular inexpensive process proxies fail to show robust discrimination under one small, conservative runner-association target, and that alternate target definitions shift the ranking substantially. Different populations, other process metrics or better semantic oracles could yield other results. Conversely, even an impressive runner-only AUROC would not establish semantic correctness or causality.
